\chapter{A guided tour of the package}
\label{ch:package}

\begin{motivation}
Theory is only useful if you can run it. This chapter maps the mathematics of the
book onto the \code{hawkes\_calibration} package: which module does what, which
experiment reproduces which figure, and how the test suite backs every numerical
claim.
\end{motivation}

\section{Module map}

\begin{center}
\renewcommand{\arraystretch}{1.2}
\begin{tabular}{@{}ll@{}}
\toprule
Module & Role (chapter)\\
\midrule
\code{eventtime/} & event-time Hawkes MLE: simulate, likelihood, estimate, lasso, covariates (Ch.~\ref{ch:hawkes})\\
\code{mbpp/core.py} & MBPP intensity/compensator, closed \& numeric (Ch.~\ref{ch:mbpp},\ref{ch:solving})\\
\code{mbpp/exogenous.py} & baseline functions incl.\ covariate link (Ch.~\ref{ch:solving},\ref{ch:baseline})\\
\code{mbpp/ic\_simulate.py} & cluster \& excitation-modulated simulators (Ch.~\ref{ch:hawkes},\ref{ch:excitation})\\
\code{mbpp/interval\_censored.py} & IC-LL/SSE, all fitters, SEs, forecasting (Ch.~\ref{ch:fitting}--\ref{ch:multitimescale})\\
\code{mbpp/gof.py} & time-rescaling KS, dispersion (Ch.~\ref{ch:identif})\\
\code{mbpp/bayes.py} & Bayesian posterior, priors, hierarchical pooling (Ch.~\ref{ch:bayes})\\
\code{operators/linear.py} & ODE/state-space, Volterra, spectral operator (Ch.~\ref{ch:solvers})\\
\code{operators/nn.py} & numpy DeepONet, amortised inference (Ch.~\ref{ch:neural})\\
\code{operators/tf.py} & TensorFlow high-dim FNO/DeepONet/SSM (Ch.~\ref{ch:neural})\\
\code{operators/neural\_solver*.py} & PINN/PINO in JAX \& TensorFlow (Ch.~\ref{ch:neural})\\
\bottomrule
\end{tabular}
\end{center}

\section{Experiments}

Each experiment is runnable (\code{python -m experiments.expN\_...}) and writes a
figure to \code{results/}:

\begin{itemize}[leftmargin=1.4em]
\item \code{exp1--4} --- event-time Hawkes: bias/variance, elicitation-matrix
recovery, covariates, sparse high-dimensional lasso.
\item \code{exp5} --- MBPP $=$ mean Hawkes intensity (Ch.~\ref{ch:mbpp}).
\item \code{exp6} --- recover $(\kappa,\theta)$ from counts across scenarios; IC-LL
vs.\ SSE (Ch.~\ref{ch:fitting}).
\item \code{exp7} --- interval-censored forecasting (ACTIVE-style).
\item \code{exp8} --- baseline covariates (Ch.~\ref{ch:baseline}).
\item \code{exp9} --- operator backends: ODE/spectral/DeepONet/amortised
(Ch.~\ref{ch:solvers},\ref{ch:neural}).
\item \code{exp10} --- excitation covariates on true Hawkes data, multi-timescale,
goodness-of-fit (Ch.~\ref{ch:excitation},\ref{ch:identif}).
\end{itemize}

\section{Reproducibility and tests}

Every numerical claim in the book is checked by the test suite
(\code{tests/test\_interval\_censored.py}, \code{test\_operators.py},
\code{test\_deep.py}, \code{test\_neural\_solver.py}) --- e.g.\ the MBPP equals the
Monte-Carlo Hawkes mean, the closed-form and numerical solvers agree, the IC-LL
recovers $\kappa$, the Volterra and LTV solvers agree where both apply, the family
residual is exactly zero, and the standard errors cover at the nominal rate on
model-correct data. The package \code{README} consolidates the practical
reference --- the capability matrix, the learned-solver (PINN/PINO) recipe, and
the TensorFlow noise-stress harness --- alongside this book.

\begin{whybox}[: tests as the book's ground truth]
Where the text asserts a number (``$\hat\kappa\approx0.60\pm0.03$'', ``residual
$=0$ to $10^{-9}$''), there is a test that produces it. Treat the suite as the
executable version of this book.
\end{whybox}

\begin{recap}
The package mirrors the book: simulation and MLE for the event-time case; the MBPP,
its solvers, and its fitters for counts; covariate and multi-timescale extensions;
the operator and neural-solver machinery; and goodness-of-fit. Experiments
reproduce the figures and tests certify the claims.
\end{recap}

\exercise{Run \code{exp6} and \code{exp10} and match each panel to a chapter result.}
\exercise{Pick one theorem from Part II and find the test that numerically
verifies it.}
